% ---------------------------------------------------------------------
\subsection{RQ5: Real LLM Proposers}\label{sec:rq5}
RQ1--RQ4 quantify over all proposals, so they say nothing about how often real models produce the situations the controller guards against. RQ5 runs the document-service workflow of Section~\ref{sec:llm-workflows} with \LlmModelsN{} open models (\LlmCallsTotal{} logged calls).

% Generated by experiments/revision_paper.py; do not edit.
\begin{table*}[!t]
\centering
\caption{RQ5: initial completion of four-service workspaces by real LLM proposers (budget $k=3$, validator feedback on retries, one feedback epoch after a barrier conflict). \emph{First-attempt acceptance}: share of first attempts committed per stage. \emph{Failed}: share of jobs that exhausted their budget (their dependents are blocked, an explicit verdict). \emph{Conflicts}: workspaces whose locally valid designs violated the global identifier convention at the barrier.}\label{tab:llm}
\footnotesize\setlength{\tabcolsep}{4pt}
\begin{tabular*}{\textwidth}{@{\extracolsep{\fill}}lrrrrrrrrr@{}}\toprule
& & \multicolumn{3}{c}{First-attempt acceptance (\%)} & & & & & \\\cmidrule(lr){3-5}
Model & Trials & Design & Code & Tests & Failed jobs (\%) & Conflicts & Published & Calls & kTokens/call\\\midrule
gemma2:9b & 6 & 0 & 0 & -- & 33 & 0 & 0 & 98 & 0.6\\
gpt-oss:20b & 6 & 12 & 100 & 67 & 3 & 0 & 4 & 104 & 1.7\\
llama3.1:8b & 6 & 0 & 4 & 0 & 29 & 0 & 0 & 138 & 0.7\\
mistral-small:24b & 6 & 0 & 58 & 83 & 1 & 0 & 5 & 113 & 0.8\\
phi4:14b & 6 & 0 & 100 & 38 & 6 & 0 & 3 & 117 & 0.8\\
qwen2.5-coder:7b & 6 & 0 & 12 & 0 & 32 & 0 & 0 & 123 & 0.7\\
qwen3-coder:30b & 6 & 0 & 92 & 25 & 4 & 0 & 4 & 122 & 0.7\\
\bottomrule\end{tabular*}
\end{table*}


\emph{Completion.} Table~\ref{tab:llm} reports \LlmTrials{} fresh four-service workspaces, six per model, each with twelve jobs. \LlmPublished{} workspaces (\LlmPublishedPct\%) reached a published snapshot. In the others, \LlmJobsFailedPct\% of all jobs exhausted their budget of $k=3$ attempts, and their dependents were blocked. Every unpublished workspace ended with an explicit verdict that named the failed jobs. No run looped, and no snapshot was published with a failed, stale or pending job. \LlmModelsNeverPublished{} of the \LlmModelsN{} models, the three with fewer than 10 billion parameters, published no workspace at all. Across all models, the most frequent implementation rejection was a violation of the stated interface (\texttt{read} must always return a pair), and the most frequent test rejection was a suite that never exercised the archived-document rule; the smaller models failed these checks again on their retries more often. The controller cannot make a weak model succeed. What it guarantees is that failure is reported, bounded and contained.

\emph{First attempts and feedback.} First-attempt acceptance of designs is close to zero, and the logs show why. \LlmDesignNameRejects{} of the \LlmDesignFirstRejects{} rejected first design attempts failed the check on the list of operation names. Of the \LlmDesignParsed{} first attempts we could parse, \LlmDesignUpdateNamed{} named the update operation \texttt{update}, whereas the validator requires \texttt{put}. That name appears in the clauses (``Updating (put) \dots'') but not in the design contract. The contract is therefore underspecified: this is a defect of our workflow, not of the models. The validator's diagnostic names the expected list, and the second attempt usually repaired it. We report the runs as executed rather than re-running with a corrected contract, because they illustrate two points. First, validator feedback recovered \LlmFeedbackRecoveredPct\% (95\% CI \LlmFeedbackRecoveredCI\%) of all rejected first attempts across stages. Second, first-attempt rates measure the clarity of a contract at least as much as the capability of a model.

\emph{Stale replies.} The stale trials are a controlled experiment, not a field observation. An implementation request is dispatched, clause R03 is flipped half a second later, and the reply is delivered after the flip. Of the 84 planned trials, 7 could not run because the design job had already failed. In the remaining \LlmStaleTrials{}, the commit guard discarded every late reply, as Theorem~\ref{thm:safety} requires. To see what accepting them would have meant, we evaluated each late reply against the validators of both questions. \LlmStaleOldValid{} replies passed the validator of the question they were asked, and \LlmStaleWrongPct\% of these (95\% CI \LlmStaleWrongCI\%) failed the validator of the current question. The other replies were invalid for both questions. Because the flip changes exactly the behavior the implementation must encode, this outcome is expected. It nonetheless quantifies a practical point: acceptance rules of the form ``the reply is valid for the question it was asked'', which is what a framework without request identities and current keys enforces, would have committed \LlmStaleOldValid{} wrong implementations.

% Generated by experiments/revision_paper.py; do not edit.
\begin{table}[!t]
\centering
\caption{RQ5: early cutoff under stochastic generation. After a reworded contract with the same meaning, how often was the regenerated value equal to the previous one? \emph{Raw}: identical response text; \emph{Value}: identical slot value (what early cutoff compares); \emph{Sem.}: same identifier type, operations and status codes; \emph{AST}: same structure-only normal form.}\label{tab:llm-cutoff}
\scriptsize\setlength{\tabcolsep}{1.6pt}
\begin{tabular*}{\columnwidth}{@{\extracolsep{\fill}}lrrrrrrrrr@{}}\toprule
& & \multicolumn{4}{c}{Design (JSON), \%} & \multicolumn{4}{c}{Implementation (code), \%}\\\cmidrule(lr){3-6}\cmidrule(lr){7-10}
Model & $T$ & $n$ & Raw & Value & Sem. & $n$ & Raw & Value & AST\\\midrule
gpt-oss:20b & 0.7 & 8 & 12 & 62 & 100 & 8 & 0 & 0 & 0\\
llama3.1:8b & 0.7 & 1 & 0 & 100 & 100 & 1 & 0 & 0 & 0\\
mistral-small:24b & 0.0 & 8 & 50 & 100 & 100 & 8 & 50 & 50 & 50\\
mistral-small:24b & 0.7 & 8 & 62 & 100 & 100 & 8 & 25 & 25 & 38\\
phi4:14b & 0.0 & 8 & 100 & 100 & 100 & 8 & 0 & 0 & 0\\
phi4:14b & 0.7 & 5 & 80 & 100 & 100 & 5 & 0 & 0 & 0\\
qwen3-coder:30b & 0.0 & 8 & 0 & 100 & 100 & 8 & 12 & 12 & 12\\
qwen3-coder:30b & 0.7 & 8 & 25 & 100 & 100 & 8 & 12 & 12 & 12\\
\bottomrule\end{tabular*}
\end{table}


\emph{Early cutoff.} Table~\ref{tab:llm-cutoff} replaces a contract by a reworded instruction with the same meaning, so a perfect proposer would regenerate an equal value. Cutoff trials need a successful initial repair, so only models that completed one appear. For the design slot, whose value is canonical JSON, the regenerated value equaled the previous one in \LlmCutoffDesignValuePct\% of trials, although the raw response text was identical in only \LlmCutoffDesignRawPct\%. Canonicalization absorbs differences in formatting and key order, which is the case for narrowly typed slots. For free-form implementations, value equality dropped to \LlmCutoffImplValuePct\% (AST equality \LlmCutoffImplAstPct\%), and most regenerations propagated to the verification job.

% Generated by experiments/revision_paper.py; do not edit.
\begin{table}[!t]
\centering
\caption{RQ5: cost of one stream of eight changes (two clause edits, three context changes of which one unselected and one equal-content, two reworded contracts) on a four-service workspace with real LLMs, after the initial repair. \emph{Publishable}: changes after which the repaired workspace passed the barrier.}\label{tab:llm-cost}
\scriptsize\setlength{\tabcolsep}{2pt}
\begin{tabular*}{\columnwidth}{@{\extracolsep{\fill}}llrrrr@{}}\toprule
Model & Strategy & Requests & kTokens & Seconds & Publishable\\\midrule
mistral-small:24b & Exact (ours) & 34 & 31.4 & 128 & 5/8\\
mistral-small:24b & No early cutoff & 44 & 41.3 & 136 & 7/8\\
mistral-small:24b & Whole workspace & 117 & 202.5 & 288 & 3/8\\
mistral-small:24b & Full regeneration & 159 & 136.5 & 419 & 7/8\\
qwen2.5-coder:7b & Exact (ours) & 50 & 40.1 & 116 & 0/8\\
qwen2.5-coder:7b & No early cutoff & 50 & 41.4 & 106 & 0/8\\
qwen2.5-coder:7b & Whole workspace & 94 & 134.4 & 116 & 0/8\\
qwen2.5-coder:7b & Full regeneration & 170 & 118.0 & 227 & 0/8\\
qwen3-coder:30b & Exact (ours) & 37 & 28.0 & 117 & 5/8\\
qwen3-coder:30b & No early cutoff & 43 & 33.3 & 143 & 5/8\\
qwen3-coder:30b & Whole workspace & 132 & 208.5 & 454 & 1/8\\
qwen3-coder:30b & Full regeneration & 154 & 107.1 & 416 & 7/8\\
\bottomrule\end{tabular*}
\end{table}


\emph{Cost and outcome.} Table~\ref{tab:llm-cost} applies one fixed stream of eight changes under four strategies with three models. Over the three models, exact keys needed \LlmCostExactReq{} requests against \LlmCostFullReq{} for full regeneration (\LlmCostRatio$\times$) and \LlmCostTokenRatio$\times$ fewer tokens. The whole-workspace strategy sent the most tokens, because every prompt carried the whole workspace, and it produced the fewest publishable snapshots. Exact keys did \emph{not}, however, dominate on outcome. For \texttt{mistral-small} they yielded a publishable snapshot after 5 of 8 changes, against 7 of 8 for full regeneration. The trace explains why. One verification job exhausted its budget at the first change, and its failed verdict persisted through the next two changes, which did not alter its observation, until a contract change opened a new episode. Full regeneration retries every job after every change and so re-samples a failed job by chance. Exact keys retain the verdict, as the semantics prescribes (Theorem~\ref{thm:termination}): Failed is terminal for an unchanged episode. This is a deliberate trade-off, not an accident. A deployment that wants more attempts can open a new episode explicitly, for example by changing the job's generation settings, which are part of its contract. It then pays for that job only, not for the whole workspace.

\textbf{Answer to RQ5.} With real models the guard was exercised as intended. Every late reply was discarded, and a large majority of those that looked valid were wrong for the current question. Validator feedback repaired most rejected attempts. Exact keys cut requests by \LlmCostRatio$\times$ against full regeneration. Early cutoff fired for canonical structured slots but rarely for free-form code. Two results qualify the picture. Smaller models often failed outright, and the controller can only make such failure explicit. And because failure verdicts are remembered, exact keys can publish less often than blind regeneration unless retries are requested explicitly. The barrier conflict that motivates the global predicate (F2) did not occur with these models: every model chose the same identifier type for all services (\LlmBarrierConflicts{} conflicts), so the evidence for the barrier remains RQ1 and the model checker.

% ---------------------------------------------------------------------
\subsection{RQ6: Validator Strength on Standard Benchmarks}\label{sec:rq6}
RQ6 asks what coordination guarantees mean for correctness when the validator is a typical test suite. It uses the \BenchTasks{} tasks of HumanEval+ and MBPP+ with \BenchModelsN{} models (\BenchCalls{} calls). The benchmark's base tests serve as the declared validator, and its plus tests serve only as hidden ground truth.

% Generated by experiments/revision_paper.py; do not edit.
\begin{table*}[!t]
\centering
\caption{RQ6: coordinated code generation on HumanEval+ (HE+) and MBPP+ with budget $k=3$ and the benchmark's base tests as the declared validator; plus tests (up to 200 inputs per task) are used only as ground truth. \emph{First correct}: first proposal passes base and plus tests (no validator, no retry). \emph{Accepted}: a value became usable. \emph{Accepted \& correct}: it also passes the plus tests. \emph{Wrong among accepted}: validator-relative correctness gap (Corollary~\ref{cor:domain}). Right block: after a reworded code contract with the same meaning, share of regenerated values that are identical as text (\emph{Raw}), identical as slot values (\emph{Value}; early cutoff fires), identical in structure (\emph{AST}), or behaviourally equal on all plus inputs (\emph{Beh.}).}\label{tab:bench}
\footnotesize\setlength{\tabcolsep}{3.5pt}
\begin{tabular*}{\textwidth}{@{\extracolsep{\fill}}llrrrrrrrrr@{}}\toprule
& & & \multicolumn{4}{c}{Correctness (\%)} & \multicolumn{4}{c}{Regeneration equal to previous value (\%)}\\\cmidrule(lr){4-7}\cmidrule(lr){8-11}
Model & Bench. & Tasks & First correct & Accepted & Accepted \& correct & Wrong among accepted & Raw & Value & AST & Beh.\\\midrule
gemma2:9b & HE+ & 164 & 61.6 & 73.2 & 65.9 & 10.0 & 61 & 63 & 75 & 97\\
gemma2:9b & MBPP+ & 378 & 61.1 & 76.2 & 64.6 & 15.3 & 62 & 71 & 75 & 96\\
llama3.1:8b & HE+ & 164 & 61.0 & 75.0 & 68.3 & 8.9 & 38 & 38 & 55 & 97\\
llama3.1:8b & MBPP+ & 378 & 64.3 & 78.3 & 68.5 & 12.5 & 11 & 13 & 59 & 91\\
mistral-small:24b & HE+ & 164 & 81.1 & 90.2 & 84.8 & 6.1 & 33 & 33 & 60 & 97\\
mistral-small:24b & MBPP+ & 378 & 67.7 & 86.2 & 71.2 & 17.5 & 6 & 6 & 11 & 95\\
phi4:14b & HE+ & 164 & 81.7 & 91.5 & 85.4 & 6.7 & 18 & 18 & 43 & 95\\
phi4:14b & MBPP+ & 378 & 71.7 & 88.6 & 73.8 & 16.7 & 6 & 6 & 20 & 91\\
qwen2.5-coder:7b & HE+ & 164 & 81.7 & 89.6 & 85.4 & 4.8 & 50 & 50 & 60 & 97\\
qwen2.5-coder:7b & MBPP+ & 378 & 71.2 & 87.3 & 74.1 & 15.2 & 21 & 21 & 42 & 93\\
\bottomrule\end{tabular*}
\end{table*}


\emph{Correctness.} Table~\ref{tab:bench} reports the results. Over all models and tasks, a value was accepted for \BenchAcceptedPct\% of tasks. \BenchAcceptedWrongPct\% of the accepted values (95\% CI \BenchAcceptedWrongCI\%) failed the plus tests: \BenchAcceptedWrongHE\% on HumanEval+ and \BenchAcceptedWrongMBPP\% on MBPP+, whose base tests are known to be weaker~\cite{evalplus}. The gap is the one Corollary~\ref{cor:domain} predicts. The controller guarantees that an accepted value passed the declared validator for the current question, and the plus tests show how much that validator leaves unchecked. Coordination neither creates this gap nor hides it. Validator-guided retry raised the share of tasks with a correct final value from \BenchFirstCorrectPct\% (first proposal, no retry) to \BenchFinalCorrectPct\%, a gain of \BenchRetryGainPts{} percentage points. Of the first proposals that passed the base tests, \BenchFirstBaseWrongPct\% still failed the plus tests.

\emph{Regeneration under a reworded contract.} After every code contract was replaced by a reworded instruction with the same meaning, \BenchRegenN{} tasks were regenerated with a usable value. The new value was identical as text in \BenchByteEqualPct\% of them, equal as a slot value in \BenchValueEqualPct\% (95\% CI \BenchValueEqualCI\%) and equal in AST structure in \BenchAstEqualPct\%. It was behaviorally equal on all plus inputs in \BenchBehEqualPct\% (95\% CI \BenchBehEqualCI\%). Early cutoff fires only on value equality, so the test job was regenerated for \BenchTestsRegenPct\% of tasks, although almost all regenerated functions behaved identically. The gap between value and behavioral equality is the saving that a coarser equality could recover, at the price of false hits (Section~\ref{sec:discussion}).

\textbf{Answer to RQ6.} On standard benchmarks, correctness of the coordinated output was bounded by the declared validator, exactly as Corollary~\ref{cor:domain} states. With the base tests as validator, about one accepted value in eight was wrong, and the rate depended on the benchmark's test strength, not on coordination. For free-form code, stochastic regeneration made exact early cutoff fire in a minority of cases, even though behavior was almost always preserved.
